\SourceLeaf{313}
Now since the denominators in II. and III. both have the same value, namely

\[
\tfrac{1}{2}\rho\left\{\sin(2A-2A')+\sin(2A'-2A'')+\sin(2A''-2A)\right\},
\]

division of these equations gives:

\[
\operatorname{tg}\psi = -\dfrac{mmh\cos A+m'm'h'\cos A'+m''m''h''\cos A''}{mmh\sin A+m'm'h'\sin A'+m''m''h''\sin A''},
\]

whence two values follow for $\psi$ and thus also for $R$.]

\bigskip
\begin{center}
\hrulefill
\end{center}
\bigskip

\Needspace{9\baselineskip}
\begin{center}
REMARK.
\end{center}

The preceding note, which is found in a handbook and presumably dates from 1816, was written very hastily, so that numerous slips in the formulas had to be corrected. It breaks off at formula II of the second method of solution, where the denominator and the factor $\Delta$ had to be supplied, something that could not be indicated in print. Two figures have been added to the text for clarity.

It may be assumed that \textsc{Gauss} found his solutions by starting from the representation of the vertices of the triangle by complex quantities, for this gives precisely the formulas in the text if each pair of equations containing the cosine and the sine of the same angle is combined into a single equation by means of the symbol $i$. That \textsc{Gauss} used the representation of points in a plane by complex quantities in geometrical problems at an early date is shown, for example, by the note on the construction for \textsc{Pothenot}'s problem printed on the following page, and by a handwritten remark on \textsc{Schumacher}'s 1810 translation of \emph{Geometrie der Stellung} by \textsc{Carnot}, which likewise concerns the geometry of the triangle and has already appeared in vol.\ IV.\ pp.\ 396 to 397.

\hfill \textsc{Stäckel}.

\bigskip
\begin{center}
\hrulefill
\end{center}
